\chapter{105}



Flink wie Vogelpicke flogen die Finger der Polizistin über die Tastatur.
Louis blieb unruhig vor ihr stehen. Während er sich wünschte, sie möge
endlos weitermachen, begann er hektisch ihre Sommersprossen zu zählen.

Bis sie eine blendend weisse Zahnreihe entblösste.

«Buongiorno, Signore.»

«Buongiorno, Signora. Ich bin Louis Vincent und habe einen Termin bei
Commissario Mantovani.»

Mit ihrem \emph{«Si»} nahm sie Louis den letzten Hoffnungsschimmer auf
die Möglichkeit, sich davon zu machen. Er wurde erwartet. Es fühlte sich
an, als wären seine Finger ferngesteuert. Vor seinen Augen schoben sie
den aufgeklappten Pass unter der Trennscheibe durch.

\sectionbreak

«War Giorgia Sala alleine am Felsen oder waren da noch weitere
Personen?»

Mantovanis Blick schien auf Louis’ Gesicht festzukleben.

«Da war niemand. Ich habe nur sie gesehen.»

Mantovani legte die Stirn in Falten und liess Louis nach wie vor nicht
aus den Augen.

«Sind sie sich absolut sicher?»

Louis’ Gehirn arbeitete beschleunigt. Sein linkes Auge fing an zu
zucken.

«Si, Commissario, wer hätte denn noch da sein sollen?»

Die Rückfrage schien ihm gelungen. Er wusste um Valentinas vergangenen
Besuch bei Giorgia und um ihre Aussagen auf der Questur. \emph{Danke,
danke, Valentina, hast du mich angerufen.} Louis wiederholte die Worte
in sich. Mit Bastiano wollte er sich keinesfalls anlegen. Ganz abgesehen
davon, dass er ihn tatsächlich nicht bei ihr gesehen hatte.

«Das kann ich Ihnen nicht sagen, Signore. Sie bleiben also dabei: Sie
sahen Giorgia Sala aus einer Distanz von acht bis zehn Metern. Und, Sie
haben keine andere Person in ihrer Nähe am Felsen gesehen?»

«Ja, Commissario, so war es, sie war ganz alleine da.»

«Fürs Protokoll bitte ich Sie den Ablauf noch einmal kurz
zusammenzufassen.»

Louis wischte sich mit dem Handrücken die Schweissperlen von der Stirn.

«Natürlich, also, ich folgte Giorgia Sala in das kleine Waldstück. Dann
sah ich sie nahe am Rande des Felsens. Sie schrie mir irgendwelche
Schimpfworte zu. Ich hatte keine Ahnung, warum und wurde wütend. Sie
wirkte eigenartig. Ich dachte, dass da etwas nicht stimmte. Dann bin ich
durchgedreht und dachte, dass ich besser verschwinde. Später hörte ich
von Valentina Ricci, dass Giorgia abgestürzt sei. Da war ich ja schon in
der Schweiz auf der Alp. Und ja, dann meinte Valentina, es sei wichtig,
dass ich diese Aussage hier bei Ihnen machen sollte. Ich würde als Zeuge
gesucht.»

«Gut, danke. Ich muss Ihnen leider die Frage noch einmal stellen,
Signore Vincent, Sie dachten nicht daran, Signore Sala Hilfe zu leisten,
sie vom Abgrund wegzuholen?»

Louis glaubte, innerlich zu verbrennen.

«Wie ich schon sagte, ich war völlig perplex, verstand nicht, was da
gerade vor sich ging und wollte nichts als weg. Jetzt im Nachhinein
mache ich mir Vorwürfe. Ich hätte ihr helfen müssen, ich, ich weiss.»

«Gut, besten Dank, Signore Vincent. Das wäre vorläufig alles. Bitte
warten Sie draussen auf das Zeugenprotokoll. Lesen Sie es aufmerksam
durch und unterschreiben Sie es. Sollten Sie Fragen dazu haben, melden
Sie sich noch einmal bei mir, okay?»

«Ja, danke Commissario, das werde ich.»

«Zum weiteren Vorgehen, Signore Vincent,» Mantovanis Stimme hörte sich
mit einem Mal deutlich energischer an, «ich werde mich die Tage bei
Ihnen melden. Halten Sie sich bitte telefonisch zu unserer Verfügung.
Eine weitere Vorladung ist wahrscheinlich.»

Louis nickte und erhob sich langsam, obgleich er nichts lieber als
augenblicklich rausgerannt wäre.

Auch Mantovani stand auf und ging um den Tisch herum mit ausgestreckter
Hand auf Louis zu.

«Arrivederci, Signore Vincent, es ist gut, dass Sie hier waren.»

\sectionbreak

Louis verliess das Gebäude über die breite Treppe auf den Innenhof der
Questur und urplötzlich tauchte der Wunsch auf. Sie sollten ihn schuldig
sprechen, ihn verurteilen, für, für irgend etwas, für «unterlassene
Hilfeleistung» vielleicht, wie Ben als Vermutung geäussert hatte. Der
Gedanke, bestraft zu werden, fühlte sich bitter und gleichzeitig entlastend an. Er
unterdrückte die Tränen und schüttelte den Kopf.

\sectionbreak

Ben fand er auf einem Barhocker eine Strasse weiter. Er war in ein reges
Gespräch mit dem Barmann verwickelt, der gerade den gefüllten Kolben in
eine Kaffeemaschine hineindrehte. Louis staunte über Bens
Italienischkenntnisse und noch weit mehr über seine Aussprache.

Er liess sich ächzend neben Ben nieder und pustete Luft aus sich heraus.
Es war geschafft. Und eine vorläufige Erlösung; es war abgelaufen, wie
geplant. Gut, hatte ihn Ben bereits auf mögliche weitere Vorladungen
vorbereitet.

Louis Anspannung fiel allmählich von ihm ab. Das lästige Nebengeräusch
jedoch blieb. Hartnäckig. Es bedrückte ihn, Mailand und seine ihm liebsten Menschen gleich
wieder zu verlassen. Der Gedanke, zu gegebener Zeit wieder hier zu sein,
half wenig.

Sie flogen mit der nächsten Maschine zurück in die Schweiz.
